%% Illustration de la parite pour l'integrale : deux aires de meme
%% signe qui se doublent (fonction paire) vs deux aires de signe
%% oppose qui s'annulent (fonction impaire). Consomme par slides.
\begin{tikzpicture}[x=0.55cm,y=1.2cm]
    \draw[-Latex] (-3.6,0) -- (3.6,0) node[right] {\small $t$};
    \draw[-Latex] (0,-0.3) -- (0,1.8);
    \fill[niceblue,opacity=0.3] (-3,0) -- plot[domain=-3:0,samples=40,variable=\t] ({\t},{1.3*exp(-\t*\t/6)}) -- (0,0) -- cycle;
    \fill[red,opacity=0.3] (0,0) -- plot[domain=0:3,samples=40,variable=\t] ({\t},{1.3*exp(-\t*\t/6)}) -- (3,0) -- cycle;
    \draw[niceblue,thick,domain=-3:3,smooth,variable=\t,samples=80]
        plot ({\t},{1.3*exp(-\t*\t/6)});
    \node[below] at (-3,0) {\small $-a$};
    \node[below] at (3,0) {\small $a$};
    \node at (0,2.1) {\small paire};
\end{tikzpicture}
\hspace{1.4cm}
\begin{tikzpicture}[x=0.55cm,y=1.2cm]
    \draw[-Latex] (-3.6,0) -- (3.6,0) node[right] {\small $t$};
    \draw[-Latex] (0,-1.6) -- (0,1.6);
    \fill[red,opacity=0.3] (-3,0) -- plot[domain=-3:0,samples=40,variable=\t] ({\t},{1.5*\t*exp(-\t*\t/6)}) -- (0,0) -- cycle;
    \fill[niceblue,opacity=0.3] (0,0) -- plot[domain=0:3,samples=40,variable=\t] ({\t},{1.5*\t*exp(-\t*\t/6)}) -- (3,0) -- cycle;
    \draw[niceblue,thick,domain=-3:3,smooth,variable=\t,samples=80]
        plot ({\t},{1.5*\t*exp(-\t*\t/6)});
    \node[below] at (-3,-0.15) {\small $-a$};
    \node[below] at (3,0) {\small $a$};
    \node at (0,1.9) {\small impaire};
\end{tikzpicture}
